\scanpage{53}
% archive sheet number: 52
Supposons $k\geq 0$, et \unsure{la|le} \unsure{dém|lemme} \unsure{fournie|prouvé} pour $k'<k$.

On sait déjà que les $p_{\deleted{\unsure{a}}i}:H^{\unsure{*|\ell}}(X)\unsure{-|\to}H^{\unsure{i|\ell}}(X)$
sont $\in\unsure{\mathcal{A}\mathrm{lg}|\mathrm{Alg}}$ pour $i<k$, donc (transp.) pour
$i>2n-k$. Soit $T_k=\mathrm{id}-\sum_{i<k}p_i-\sum_{j>2n-k}p_j$, donc
$T_k\in\unsure{\mathcal{A}|A}$, plus \unsure{gén.} $T_{k'}\in\mathcal{A}$ si $k'\unsure{\leq|<}k$. Alors
$\unsure{\xi|\zeta}^{n-k}T_k$ est nul sur $H^{\unsure{i|j}}$, $j\neq k$, et induit
$H^k\xrightarrow{\overset{\xi^{n-k}}{\sim}}H^{2n-k}$, donc $u_{\unsure{k|n}}\xi^{n-k}T_k=p_k$.

D'autre part, si $i+2j=k$\inserted{, $j>0$}, la projection sur
$\xi^j\unsure{P|p}^{i}$ est aussi $\in\mathcal{A}$, comme on voit \deleted{\unsure{ainsi|aussitôt}}
par récurrence ascendante sur $i$ (donc descendante \unsure{sur|en}
$j$) ainsi : si $i<0$ rien à prouver. Si c'est
fait pour les $i'<i$, on forme
$\xi^{n-k+j}T_k$, \deleted{\unsure{$H^j$}} qui induit \unsure{zéro|0|nul} sur $H^\alpha$ pour $\alpha\neq k$,
et induit $H^{k}\deleted{\unsure{\xi}}\xrightarrow{\xi^{\unsure{n-k-j|n-k+j}}}H^{2n-k+2j}$, \deleted{\unsure{qui|facteurs}}\inserted{qui est nul}
sur les facteurs $\xi^{j'}P^{i'}$ pour $j'<j$, \inserted{de somme $N$,} et sur
les facteurs $\xi^{j'}P^{i'}$ avec $j'\geq j$ (i.e. $i'\leq i$),
\deleted{\unsure{coïncide avec} $1$} dont la somme n'est autre
que $\xi^j H^{k-2j}$, son composé avec l'\unsure{isom}
$H^{k-2j}\xrightarrow{\xi^{j}}\xi^j H^{k-2j}$ est $\xi^{\unsure{k|n}-(k-2j)}$. Donc l'\unsure{hom|homom.}
$u_{k-2j}\xi^{n-k+j}T_k$\unsure{:|,|;}, nul sur $N$, est sur $\xi^j H^{k-2j}$
l'inverse de $H^{k-2j}\xrightarrow{\sim}\xi^j H^{k-2j}$. Donc

